% !TeX root = ../../main.tex

\section{Conceptos básicos de Big Data}

\begin{enumerate}

\item \textbf{¿Qué es Big Data?}

Big Data es un conjunto de datos cuyo volumen, velocidad o variedad supera la capacidad de las herramientas tradicionales para almacenarlos, procesarlos y analizarlos de manera eficiente .

\item \textbf{Históricamente, ¿cómo han evolucionado las fuentes o generadores de datos?}

La generación de datos ha evolucionado de la siguiente manera:

\begin{itemize}
    \item Antes de 1980, los datos eran generados principalmente por dispositivos.
    \item Entre 1980 y 2000, los empleados generaban datos al utilizar sistemas informáticos.
    \item Desde el año 2000, las personas comenzaron a producir datos mediante correos electrónicos, redes sociales y páginas web.
    \item Después de 2005, el hardware, el software, las aplicaciones y los dispositivos del Internet de las cosas comenzaron a generar grandes cantidades de registros automáticamente.
\end{itemize}

\item \textbf{¿Qué y cuánto se genera cada 60 segundos en el mundo digital? Datos de 2023.}

Según el informe \textit{Data Never Sleeps 11.0}, durante cada minuto de 2023 se realizaron, entre otras, las siguientes actividades:

\begin{itemize}
    \item 6.3 millones de búsquedas en Google.
    \item 41.6 millones de mensajes enviados por WhatsApp.
    \item 241 millones de correos electrónicos enviados.
    \item 360\,000 publicaciones realizadas en X.
    \item 694\,000 reels enviados por mensaje directo en Instagram.
    \item 4 millones de publicaciones recibieron ``Me gusta'' en Facebook.
    \item 6\,944 solicitudes fueron enviadas a ChatGPT.
    \item Se consumieron 43 años de contenido en servicios de transmisión.
    \item Se reprodujeron 24\,000 horas de música en Spotify.
    \item Se realizaron más de \$455\,000 en compras en Amazon.
    \item Se enviaron \$463\,768 mediante Venmo.
    \item Se realizaron \$122\,785 en pedidos de DoorDash.
    \item Se produjeron 30 ataques de denegación de servicio distribuido (DDoS).
\end{itemize}

\item \textbf{¿Por qué es necesario almacenar y procesar datos en Big Data?}

Es necesario almacenar datos históricos para aprender de ellos, encontrar patrones y tomar mejores decisiones. Su procesamiento permite descubrir información útil, resolver problemas, reducir riesgos y crear nuevas oportunidades.

\item \textbf{Mencione ejemplos de aplicaciones de Big Data.}

Algunos ejemplos son:

\begin{itemize}
    \item Recomendaciones de productos en tiendas en línea.
    \item Detección de fraudes bancarios.
    \item Análisis de redes sociales y opiniones.
    \item Optimización del tráfico y las rutas de transporte.
    \item Análisis de expedientes médicos y genomas.
    \item Predicción de enfermedades y desastres naturales.
    \item Análisis del mercado de valores.
    \item Seguridad y prevención de delitos.
\end{itemize}

\item \textbf{Mencione las principales características de Big Data.}

Las principales características se conocen como las ``V'' de Big Data:

\begin{itemize}
    \item \textbf{Volumen:} gran cantidad de datos.
    \item \textbf{Velocidad:} rapidez con la que los datos se generan y procesan.
    \item \textbf{Variedad:} existencia de diferentes formatos y fuentes.
    \item \textbf{Valor:} utilidad que puede obtenerse de los datos.
    \item \textbf{Veracidad:} nivel de exactitud y confiabilidad.
    \item \textbf{Variabilidad:} cambios en el comportamiento y significado de los datos.
    \item \textbf{Volatilidad:} tiempo durante el cual los datos siguen siendo útiles.
    \item \textbf{Complejidad:} relaciones existentes entre grandes cantidades de variables.
\end{itemize}

\item \textbf{Describa las principales unidades de medida de datos.}

El bit es la unidad mínima y puede tener el valor 0 o 1. Un byte está formado por 8 bits. Las unidades principales son:

\begin{itemize}
    \item Kilobyte (KB): 1\,024 bytes.
    \item Megabyte (MB): 1\,024 KB.
    \item Gigabyte (GB): 1\,024 MB.
    \item Terabyte (TB): 1\,024 GB.
    \item Petabyte (PB): 1\,024 TB.
    \item Exabyte (EB): 1\,024 PB.
    \item Zettabyte (ZB): 1\,024 EB.
    \item Yottabyte (YB): 1\,024 ZB.
\end{itemize}

\item \textbf{Explique los diferentes tipos de datos.}

\begin{itemize}
    \item \textbf{Estructurados:} tienen una organización y un esquema definido, como las tablas de una base de datos.
    \item \textbf{Semiestructurados:} poseen cierta organización, pero no una estructura rígida. Algunos ejemplos son XML, HTML, JSON, correos electrónicos y registros de aplicaciones.
    \item \textbf{No estructurados:} no tienen un esquema fijo. Por ejemplo, imágenes, videos, audios y documentos de texto.
\end{itemize}

Estos tres tipos pueden formar parte de un sistema de Big Data

\item \textbf{¿Cuáles son las diferencias entre datos estructurados y no estructurados?}

Los datos estructurados tienen campos y formatos definidos, por lo que pueden almacenarse fácilmente en tablas y consultarse mediante SQL. Los datos no estructurados no siguen un esquema fijo, presentan formatos variados y requieren herramientas especiales para almacenarlos y analizarlos.

\item \textbf{¿Cuáles son los tipos de sistemas de Big Data?}

Existen dos tipos principales

\begin{itemize}
    \item \textbf{Sistemas operacionales:} procesan actividades diarias y proporcionan respuestas rápidas o en tiempo real.
    \item \textbf{Sistemas analíticos:} procesan grandes conjuntos de datos para encontrar patrones y apoyar la toma de decisiones.
\end{itemize}

\item \textbf{Defina los siguientes conceptos.}

\begin{itemize}
    \item[a)] \textbf{Transacción:} conjunto de operaciones relacionadas que deben realizarse correctamente y en un orden determinado. Por ejemplo, retirar dinero de una cuenta y depositarlo en otra.

    \item[b)] \textbf{Base de datos transaccional:} base de datos que permite realizar transacciones de forma segura y consistente. Se utiliza, por ejemplo, en bancos y sistemas financieros.

    \item[c)] \textbf{Base de datos operacional:} base de datos diseñada para las operaciones cotidianas, como insertar, consultar, actualizar y eliminar información. Debe responder rápidamente y puede utilizar tecnologías NoSQL.
\end{itemize}


\item \textbf{¿Cuáles son las diferencias entre los sistemas de Big Data operacional y analítico?}

\begin{itemize}
    \item El sistema operacional responde en milisegundos; el analítico puede tardar minutos u horas.
    \item El operacional atiende a muchos usuarios al mismo tiempo; el analítico suele ser utilizado por pocos especialistas.
    \item El operacional realiza lecturas, escrituras y actualizaciones frecuentes; el analítico trabaja principalmente con cargas y lecturas masivas.
    \item El operacional utiliza comúnmente bases NoSQL; el analítico emplea almacenes de datos, Hadoop o Spark.
    \item El operacional apoya actividades diarias; el analítico ayuda a encontrar patrones y tomar decisiones.
\end{itemize}

\item \textbf{¿Qué es el escalamiento horizontal?}

Consiste en agregar más computadoras o nodos a un grupo para distribuir los datos y procesarlos en paralelo. Los equipos trabajan juntos como un solo sistema lógico.

\item \textbf{¿Qué es el escalamiento vertical?}

Consiste en aumentar los recursos de una sola computadora, por ejemplo, agregando memoria, almacenamiento o procesadores más potentes.

\item \textbf{Mencione las diferencias entre el escalamiento vertical y horizontal.}

\begin{itemize}
    \item El escalamiento horizontal agrega más computadoras; el vertical mejora una sola computadora.
    \item El horizontal puede crecer agregando nodos; el vertical está limitado por la capacidad máxima del hardware.
    \item En el horizontal, la falla de un nodo no necesariamente detiene todo el sistema; en el vertical, la falla del servidor puede detener el servicio.
    \item El horizontal utiliza hardware común y suele ser más económico; el vertical requiere equipos más potentes y costosos.
    \item El horizontal necesita software distribuido más complejo; el vertical es más sencillo de administrar.
    \item El horizontal puede ampliarse sin detener el sistema; el vertical normalmente requiere apagar o reiniciar el servidor.
\end{itemize}

\end{enumerate}